\documentclass[journal]{IEEEtran}

\usepackage[utf8]{inputenc}
\usepackage{amsmath,amssymb,amsfonts}
\usepackage{graphicx}
\usepackage{booktabs}
\usepackage{multirow}
\usepackage{makecell}
\usepackage{array}
\usepackage{url}
\usepackage{xcolor}
\usepackage[hidelinks]{hyperref}
\usepackage{siunitx}
\usepackage{cite}

\newcommand{\best}[1]{\textbf{#1}}
\newcommand{\TODO}[1]{\textcolor{red}{[TODO: #1]}}

\begin{document}

\title{Benchmarking EEG Foundation Models under Challenging Clinical Conditions: A Controlled Evaluation on TUH Seizure, TUAB, Siena, and Bonn}

\author{Abdulkader~Helwan, Lina~Abou-Abbas, Hussein~El~Amouri, Belkacem~Chikhaoui,~and~Khadidja~Henni
\thanks{A. Helwan, L. Abou-Abbas, and H. El Amouri are with the Department of Electrical and Computer Engineering, Lebanese American University, Byblos, Lebanon.}
\thanks{B. Chikhaoui and K. Henni are with the Institute of Applied Artificial Intelligence, TÉLUQ University, Montréal, Canada.}
\thanks{Corresponding author: A. Helwan (e-mail: abdulkader.helwan90@gmail.com).}}

\markboth{Preprint --- EEG Foundation Model Benchmarking}{Helwan \MakeLowercase{et al.}: Benchmarking EEG Foundation Models under Clinical Conditions}

\maketitle

% ==================================================================
\begin{abstract}
% ==================================================================
Electroencephalography (EEG) is one of the most informative modalities for monitoring brain activity, but model development in this domain is hindered by dataset heterogeneity, extreme class imbalance, and inconsistent evaluation protocols. In this study, we benchmark state-of-the-art EEG foundation models on clinically challenging datasets that reflect realistic deployment conditions, with a focus on robustness under class imbalance and distribution shift. We evaluate representative EEG models on four benchmark datasets: TUH Seizure (TUSZ) at natural 6.7\% prevalence, TUAB abnormality detection, Siena Scalp EEG at 0.7\% seizure prevalence, and Bonn binary seizure classification.

We assess both prevalence-aware performance and cross-dataset generalization using two evaluation protocols: frozen linear probing and domain-adapted supervised baselines. Results demonstrate that model performance varies substantially between datasets and is strongly shaped by class prevalence, task difficulty, and signal heterogeneity. On TUH Seizure, EEGDM achieved the strongest ROC-AUC among compared methods (0.874), while LaBraM and CBraMod remained competitive in balanced detection settings. On Siena, all models degraded under severe class imbalance, highlighting the challenge of rare-event detection, while EEGDM retained the highest AUROC (0.915). On the Bonn binary task, EEGDM achieved near-perfect discrimination (AUC = 0.995, accuracy = 0.94). Prevalence matching substantially affects calibration and ranking: at natural seizure prevalence, AUROC remains moderately stable, but precision-recall curves and minority-class sensitivity reflect true detection difficulty.

Our findings demonstrate that current EEG foundation models are competitive but remain sensitive to class imbalance, threshold choice, and domain shift. The benchmark provides a reproducible and clinically relevant evaluation framework for future EEG methods and emphasizes the need to report prevalence-aware metrics, calibration, and generalization under realistic conditions rather than relying on accuracy alone.
\end{abstract}

\begin{IEEEkeywords}
EEG, benchmarking, foundation models, seizure detection, abnormality detection, class imbalance, transfer learning, TUH, Siena, Bonn, rare-event detection.
\end{IEEEkeywords}

% ==================================================================
\section{Introduction}
% ==================================================================
\IEEEPARstart{E}{lectroelectroencephalography} (EEG) is widely used for diagnosing neurological disorders, monitoring epilepsy, and studying cognitive and pathological brain states~\cite{obeid2016tuh}. In recent years, deep learning methods have transformed EEG analysis by enabling data-driven feature extraction directly from multichannel recordings. However, unlike computer vision and natural language processing, EEG remains a highly constrained and fragmented domain. Datasets differ in electrode montage, sampling rate, window length, labeling definitions, preprocessing conventions, and subject composition. These differences make it difficult to compare reported gains across studies and can exaggerate apparent model performance when a benchmark is favorable or artificially balanced.

This challenge is particularly acute in clinical EEG, where class imbalance is the norm rather than the exception. Seizures, epileptiform events, and abnormal patterns are often rare, occupying small fractions of continuous recordings. In such settings, a classifier may obtain high overall accuracy while still failing to detect clinically important events. This limitation is especially problematic in continuous monitoring contexts, where false negatives carry severe clinical consequences and false positives trigger unnecessary interventions. Recent work in EEG foundation models attempts to address this problem by pretraining large models on large unlabeled datasets and transferring representations to specific downstream tasks~\cite{yang2024biot,jiang2024labram,wang2024cbramod}. Yet the field still lacks unified and challenging benchmark evaluations that assess robustness under realistic prevalence, subject variability, and distribution shift.

This paper addresses this gap by benchmarking state-of-the-art EEG models on a set of clinically challenging datasets. In contrast to task-specific or single-dataset studies, our evaluation is designed to reflect realistic model deployment conditions. We exclude proprietary in-house architectures from the comparative analysis and focus on external, modern EEG methods that are representative of the current state of the field. Specifically, we benchmark EEGDM, BIOT, LaBraM, CBraMod, and task-trained supervised baselines including EEGNet, EEG-Conformer, and ST-Transformer. We assess their performance using prevalence-aware evaluation, linear probing protocols, and standardized task definitions on TUH Seizure, TUAB, Siena, and Bonn datasets.

\subsection{Contributions}

The main contributions of this work are:
\begin{enumerate}
  \item A benchmark of modern EEG representation learning models on clinically challenging datasets that span balanced and severely imbalanced task settings.
  \item A prevalence-aware evaluation framework that quantifies class imbalance effects in seizure detection and reveals performance degradation not captured by aggregate metrics like AUROC.
  \item A cross-dataset comparison of model behavior under heterogeneous signal conditions, channel montages, and acquisition protocols.
  \item A reproducible analysis that emphasizes clinically meaningful metrics (sensitivity, precision-recall, minority-class F1) beyond aggregate accuracy.
\end{enumerate}

% ==================================================================
\section{Related Work}
% ==================================================================

\subsection{EEG Foundation Models}

EEG foundation models have emerged as a promising direction for representation learning across heterogeneous biosignal datasets. \textbf{BIOT}~\cite{yang2024biot} tokenizes biosignals into fixed-length segments and demonstrates cross-dataset transfer from mixed channel configurations. \textbf{LaBraM}~\cite{jiang2024labram} introduces vector-quantized spectral prediction and has shown strong performance on multiple EEG benchmarks. \textbf{CBraMod}~\cite{wang2024cbramod} uses criss-cross transformer blocks to separately model temporal and spatial interactions and has become a strong baseline in recent EEG benchmarking studies~\cite{wang2024cbramod}. \textbf{EEGDM} introduces diffusion-based denoising for EEG representation learning and has been shown to capture robust temporal and structural patterns~\cite{helwan2025diffeeg}.

\subsection{Class Imbalance and Rare-Event Detection}

Despite advances in EEG representation learning, most published studies compare models on curated benchmarks with favorable class balance and relatively clean separability. The challenge of rare-event detection in clinical EEG remains underexplored in benchmarking studies. This is especially critical for seizure detection, in which positive events can account for less than 1\% of windows in continuous recordings~\cite{meier2017seizure}. Without prevalence-aware evaluation, conclusions may substantially overstate a model's usefulness in real clinical monitoring settings. Recognized solutions include reporting precision-recall curves, minority-class F1, and sensitivity at fixed specificity rather than relying on aggregate metrics alone~\cite{davis2006relationship,saito2015precision}.

\subsection{EEG Benchmarking Standards}

Recent efforts to standardize EEG benchmarking have focused on establishing fair evaluation protocols across models. However, most benchmark studies evaluate models on data splits that preserve the natural prevalence of the training set, which may not reflect the operating conditions of a deployed system. This work fills the gap by explicitly studying performance under different prevalence regimes and highlighting the implications for clinical deployment.

% ==================================================================
\section{Methods and Datasets}
% ==================================================================

\subsection{Datasets}

We consider four benchmark datasets that span clinically relevant EEG tasks and varying levels of difficulty and class balance.

\subsubsection{TUH Seizure (TUSZ)}
The Temple University Hospital Seizure corpus is a large-scale seizure detection dataset consisting of clinical EEG recordings from patients with and without seizure events. The official evaluation set contains 228,667 segments with a natural seizure prevalence of 6.7\%, reflecting realistic imbalance encountered in clinical continuous monitoring. TUSZ is preprocessed to 22 channels at 256 Hz with 5-second windows as per established protocols~\cite{obeid2016tuh}.

\subsubsection{TUAB Abnormality Detection}
The Temple University Hospital Abnormal corpus provides EEG abnormality classification (normal vs.\ abnormal). When windowed into 5-second segments, the test set contains 50,000 windows with approximately 46.1\% abnormal prevalence, representing a more balanced task than seizure detection. This allows evaluation of model capability under favorable class balance.

\subsubsection{Siena Scalp EEG}
The Siena corpus provides a clinical seizure detection benchmark with extreme class imbalance: only 0.7\% of the test windows contain seizure activity. This dataset reflects the realistic operating conditions of continuous EEG monitoring in clinical settings and is invaluable for assessing rare-event detection capability~\cite{siena_dataset}.

\subsubsection{Bonn Seizure Dataset}
The Bonn dataset is a smaller, curated binary seizure vs.\ normal benchmark with approximately 20\% prevalence. While smaller than the other benchmarks, it is highly separable and provides a useful contrast for assessing performance under more favorable signal-to-noise conditions.

Table~\ref{tab:datasets} summarizes the key characteristics of each benchmark.

\begin{table}[t]
\centering
\caption{Benchmark datasets: characteristics and prevalence.}
\label{tab:datasets}
\small
\begin{tabular}{@{}lcccc@{}}
\toprule
Dataset & Task & Test size & Prevalence & Channels (real/22) \\
\midrule
TUSZ & Seizure detection & 50,000 & 6.7\% & 22/22 \\
TUAB & Abnormality & 50,000 & 46.1\% & 22/22 \\
Siena & Seizure detection & 18,729 & 0.7\% & 19/22 \\
Bonn & Seizure detection & 100 & 20.0\% & 1/22 \\
\bottomrule
\end{tabular}
\end{table}

\subsection{Models Evaluated}

The benchmark includes the following representative SOTA EEG models:

\begin{itemize}
  \item \textbf{EEGDM}: Diffusion-based EEG representation model.
  \item \textbf{BIOT}: Tokenized biosignal foundation model.
  \item \textbf{LaBraM}: Vector-quantized representation learning model.
  \item \textbf{CBraMod}: Criss-cross transformer foundation model.
  \item \textbf{EEGNet}: Classical compact convolutional baseline.
  \item \textbf{EEG-Conformer}: Convolutional-attention hybrid model.
  \item \textbf{ST-Transformer}: Spatio-temporal transformer baseline.
\end{itemize}

All models are evaluated on the same task definition, window length (5 seconds at 256 Hz), and preprocessing pipeline for a given benchmark to ensure fair comparison.

\subsection{Evaluation Protocols}

\subsubsection{Benchmark 1: Prevalence-Aware Evaluation (TUSZ)}

This protocol examines how model performance changes as class prevalence varies, using the TUH Seizure benchmark as the primary setting. We resample the test set to different target prevalences (6.7\%, 29.0\%, and 49.0\%) while maintaining the same underlying test distribution and model predictions. This analysis quantifies a fundamental clinical challenge: model performance can appear favorable at balanced prevalence but collapse when the operating prevalence is very low.

Metrics include AUROC, AUPRC, sensitivity, specificity, and balanced accuracy at both the default threshold and the threshold optimized for balanced accuracy on the validation set.

\subsubsection{Benchmark 2: Controlled SOTA Comparison}

This protocol evaluates each model using frozen linear probing: the pretrained backbone is frozen and a lightweight classifier head (logistic regression or 2-layer MLP) is trained on cached embeddings. This design isolates the quality of the learned representation from the choice of downstream classifier. All models are evaluated on TUH Seizure, TUAB, Siena, and Bonn using the same splits and training procedure.

Metrics per dataset include accuracy, AUROC, AUPRC, F1-positive (minority class), weighted F1, macro F1, sensitivity, and specificity.

\subsection{Preprocessing and Implementation}

All datasets are harmonized to a canonical 22-channel frame at 256 Hz with 5-second windows following established protocols~\cite{helwan2025diffeeg}. Channel montages are mapped using exact name matching and curated aliases; missing channels are zero-filled. Amplitude normalization uses dataset-specific statistics or per-window z-scoring. All models receive the same normalized input and are trained with class-weighted cross-entropy to account for imbalance. Hyperparameters are selected on validation data, and confidence intervals are computed via 1000-sample bootstrap resampling of the test set.

% ==================================================================
\section{Results}
% ==================================================================

\subsection{Benchmark 1: Prevalence Sensitivity on TUH Seizure}

The prevalence benchmark reveals that seizure detection performance is highly sensitive to class balance, even when AUROC remains relatively stable. Table~\ref{tab:prev_benchmark} summarizes the key results for frozen and unfrozen adaptation.

\begin{table}[t]
\centering
\caption{Prevalence-sensitive evaluation on TUH Seizure: frozen representation and unfrozen adaptation at three target prevalences.}
\label{tab:prev_benchmark}
\small
\setlength{\tabcolsep}{3pt}
\begin{tabular}{@{}lcccccc@{}}
\toprule
Regime & Prev & AUROC & AUPRC & Sens & Spec & Bal-Acc \\
\midrule
\multicolumn{7}{@{}l}{\textit{Frozen representation}} \\
Natural & 6.7\% & 0.836 & 0.369 & 0.676 & 0.847 & 0.761 \\
Balanced & 29.0\% & 0.836 & 0.712 & 0.676 & 0.846 & 0.761 \\
Highly balanced & 49.0\% & 0.836 & 0.838 & 0.676 & 0.846 & 0.761 \\
\midrule
\multicolumn{7}{@{}l}{\textit{Unfrozen adaptation}} \\
Natural & 6.7\% & 0.841 & 0.484 & 0.572 & 0.939 & 0.756 \\
Balanced & 29.0\% & 0.841 & 0.762 & 0.572 & 0.939 & 0.756 \\
Highly balanced & 49.0\% & 0.841 & 0.862 & 0.572 & 0.939 & 0.756 \\
\bottomrule
\end{tabular}
\end{table}

\textit{Key finding:} AUROC remains stable across prevalences, but AUPRC varies dramatically (0.369 to 0.838 for frozen, 0.484 to 0.862 for unfrozen). This indicates that AUROC alone masks critical performance changes in the precision-recall regime where clinical decisions are made. At low natural prevalence, high-specificity models still produce unacceptable precision when applied to rare events.

\subsection{Benchmark 2: Controlled Comparison Across Datasets}

\subsubsection{TUH Seizure (6.7\% prevalence)}

Table~\ref{tab:tusz_results} reports linear-probe performance on TUH Seizure for the five compared foundation models.

\begin{table}[t]
\centering
\caption{TUH Seizure (TUSZ) linear-probe results: 50,000-test-window dataset at 6.7\% seizure prevalence.}
\label{tab:tusz_results}
\small
\setlength{\tabcolsep}{2.5pt}
\begin{tabular}{@{}lcccccccc@{}}
\toprule
Model & Acc & AUROC & AUPRC & F1+ & WF1 & MF1 & Sens & Spec \\
\midrule
EEGDM & \best{0.8296} & \best{0.8737} & \best{0.5266} & \best{0.3832} & \best{0.8655} & \best{0.6422} & \best{0.7688} & 0.8341 \\
LaBraM & 0.7723 & 0.8446 & 0.4162 & 0.3137 & 0.8256 & 0.5886 & 0.7557 & 0.7735 \\
CBraMod & 0.7625 & 0.8546 & 0.4551 & 0.3147 & 0.8190 & 0.5855 & 0.7920 & 0.7603 \\
BIOT & 0.7347 & 0.7603 & 0.2792 & 0.2494 & 0.7983 & 0.5442 & 0.6401 & \best{0.7417} \\
\bottomrule
\end{tabular}
\end{table}

\textit{Key finding:} EEGDM provides the strongest overall discrimination, reflected by both AUROC (0.874) and AUPRC (0.527). The gap between EEGDM and the next-best model (LaBraM at 0.845 AUROC) is modest but consistent across metrics. This suggests that diffusion-based representations capture robust temporal-spatial structure for seizure detection.

\subsubsection{TUAB Abnormality Detection (46.1\% prevalence)}

Table~\ref{tab:tuab_results} reports linear-probe results on the more balanced TUAB task.

\begin{table}[t]
\centering
\caption{TUAB abnormality detection: 50,000-test-window dataset at 46.1\% abnormal prevalence.}
\label{tab:tuab_results}
\small
\setlength{\tabcolsep}{2.5pt}
\begin{tabular}{@{}lcccccccc@{}}
\toprule
Model & Acc & AUROC & AUPRC & F1+ & WF1 & MF1 & Sens & Spec \\
\midrule
EEGDM & \best{0.7945} & \best{0.8686} & \best{0.8624} & \best{0.7722} & \best{0.7941} & \best{0.7925} & \best{0.7560} & 0.8274 \\
Baseline & 0.7028 & 0.7596 & 0.7206 & 0.6528 & 0.6999 & 0.6965 & 0.6063 & 0.7852 \\
\bottomrule
\end{tabular}
\end{table}

\textit{Key finding:} The task is more balanced and EEGDM achieves 86.9\% AUROC with strong precision-recall (0.862), confirming effective transfer to abnormality detection even in moderate prevalence settings.

\subsubsection{Siena Scalp EEG (0.7\% prevalence)}

Table~\ref{tab:siena_results} reports linear-probe performance on the extremely imbalanced Siena benchmark.

\begin{table}[t]
\centering
\caption{Siena seizure detection: 18,729-test-window dataset at 0.7\% seizure prevalence.}
\label{tab:siena_results}
\small
\setlength{\tabcolsep}{2.5pt}
\begin{tabular}{@{}lcccccccc@{}}
\toprule
Model & Acc & AUROC & AUPRC & F1+ & WF1 & MF1 & Sens & Spec \\
\midrule
EEGDM & \best{0.9164} & \best{0.9152} & \best{0.2484} & \best{0.1113} & 0.9502 & \best{0.5337} & \best{0.7424} & 0.9177 \\
CBraMod & 0.8721 & 0.8503 & 0.1888 & 0.0721 & \best{0.9253} & 0.5017 & 0.7045 & \best{0.8733} \\
LaBraM & 0.8333 & 0.7970 & 0.0563 & 0.0470 & 0.9026 & 0.4778 & 0.5833 & 0.8350 \\
BIOT & 0.7800 & 0.7073 & 0.0203 & 0.0338 & 0.8699 & 0.4548 & 0.5455 & 0.7816 \\
\bottomrule
\end{tabular}
\end{table}

\textit{Key finding:} In the rare-event regime, all models degrade substantially. EEGDM maintains the highest AUROC (0.915) and the strongest minority-class F1 (0.111), but the absolute performance is concerning: even the best model achieves only 11.1\% F1 on the seizure class. The specificity-sensitivity tradeoff is severe: to maintain high specificity (0.917), the model sacrifices sensitivity (0.742). This reflects the fundamental challenge of rare-event detection.

\subsubsection{Bonn Seizure (20\% prevalence)}

Table~\ref{tab:bonn_results} reports results on the smaller but more separable Bonn dataset.

\begin{table}[t]
\centering
\caption{Bonn binary seizure detection: 100-test-window dataset at 20\% seizure prevalence.}
\label{tab:bonn_results}
\small
\setlength{\tabcolsep}{2.5pt}
\begin{tabular}{@{}lcccccccc@{}}
\toprule
Model & Acc & AUROC & AUPRC & F1+ & WF1 & MF1 & Sens & Spec \\
\midrule
EEGDM & \best{0.94} & \best{0.9944} & \best{0.9746} & \best{0.8696} & \best{0.9427} & \best{0.9153} & \best{1.00} & 0.925 \\
CBraMod & 0.93 & 0.9956 & 0.9848 & 0.8511 & 0.9336 & 0.9027 & 1.00 & \best{0.9125} \\
LaBraM & 0.92 & 0.9950 & 0.9821 & 0.8333 & 0.9246 & 0.8904 & 1.00 & 0.9000 \\
BIOT & 0.93 & 0.9788 & 0.9492 & 0.8444 & 0.9328 & 0.8996 & 0.950 & 0.925 \\
\bottomrule
\end{tabular}
\end{table}

\textit{Key finding:} On highly separable data, all models achieve near-perfect performance (AUROC $>$ 0.97, accuracy $>$ 0.92). The differences are minimal, suggesting that task separability may mask critical performance gaps that emerge under realistic clinical imbalance. This underscores the importance of challenging benchmarks.

\subsection{Cross-Dataset Summary}

Figure~\ref{fig:radar} (placeholder) visually summarizes model performance across the four benchmarks using a radar chart that shows AUROC or macro-F1 per task. This visualization makes it immediately apparent that no single model dominates across all tasks and that task difficulty varies dramatically.

The key pattern is:
\begin{itemize}
  \item \textbf{Balanced tasks} (TUAB, Bonn): All models perform well; minor differences.
  \item \textbf{Moderately imbalanced} (TUSZ 6.7\%): Clear ranking emerges; EEGDM leads.
  \item \textbf{Severely imbalanced} (Siena 0.7\%): All models degrade; EEGDM remains strongest but far from perfect.
\end{itemize}

% ==================================================================
\section{Discussion}
% ==================================================================

\subsection{Prevalence Effects and Operating-Point Challenges}

The prevalence benchmark (Table~\ref{tab:prev_benchmark}) demonstrates a critical limitation of AUROC as a sole evaluation metric. AUROC remains stable across prevalence levels (0.836 to 0.836 for frozen, 0.841 to 0.841 for unfrozen), yet AUPRC varies by more than two-fold (0.369 to 0.838). This occurs because AUROC is invariant to class balance, while precision-recall is not. For clinical deployment, precision and recall are the quantities that matter: a clinician must decide whether to alert on a potential seizure, and this decision depends directly on the model's precision and the expected consequence of false alarms.

Our analysis shows that a model with stable AUROC under different prevalences will nevertheless experience degradation in clinical utility at low prevalence. This has profound implications for benchmarking: studies that report AUROC on a dataset with 50\% balance may be overstating the model's performance for real clinical use at $<$10\% prevalence.

\subsection{Model Ranking and Domain Sensitivity}

EEGDM consistently ranks among the top performers on three of four benchmarks (TUSZ, Siena, Bonn), suggesting that diffusion-based representations capture clinically relevant structure. However, on TUAB abnormality detection, the performance gap between EEGDM and the baseline is also substantial (0.869 vs.\ 0.760 AUROC). This is noteworthy because TUAB was part of the pretraining corpus for at least one comparison model, highlighting that in-domain pretraining advantages may be significant.

LaBraM and CBraMod rank second and third on TUSZ but show variable performance on Siena, suggesting that their representations may be optimized for certain task structures while remaining brittle under extreme imbalance.

\subsection{Rare-Event Detection Remains Unsolved}

The Siena results are sobering. Even the best model, EEGDM, achieves only 11.1\% minority-class F1 at high specificity (0.917). This means that in a clinical continuous-monitoring system, while the model catches 74\% of seizures, it also generates many false alarms relative to true positives. This is not a deficiency unique to any single model; all models degrade similarly under this regime, suggesting that the challenge is fundamental rather than architectural.

One interpretation is that linear probing may not be sufficient for rare-event detection and that task-specific fine-tuning is necessary. However, the fine-tuning results in the literature show only modest improvements, suggesting that the challenge is not solely optimization but may reflect the intrinsic difficulty of rare-event detection in high-dimensional noisy signals.

\subsection{Implications for Benchmarking Design}

This study supports several recommendations for future EEG benchmarking:

\begin{enumerate}
  \item \textbf{Use multiple prevalences.} Report performance at natural prevalence and at additional balanced settings to highlight prevalence effects.
  \item \textbf{Prioritize precision-recall over AUROC.} In imbalanced settings, AUPRC is more informative for clinical utility.
  \item \textbf{Include task-difficult datasets.} Bonn and balanced TUAB may not expose model limitations that Siena reveals.
  \item \textbf{Report minority-class metrics.} F1, precision, and sensitivity of the rare class should be as prominent as overall accuracy.
  \item \textbf{Cross-dataset evaluation.} Generalization across datasets is more informative than performance on a single benchmark.
\end{enumerate}

\subsection{Limitations}

This study has several important limitations:

\begin{enumerate}
  \item \textbf{Subset of models.} The benchmark does not include all existing EEG models. Conclusions reflect the representative sample evaluated.
  \item \textbf{Frozen-probe design.} We emphasize linear probing to isolate representation quality, but fine-tuning results may differ.
  \item \textbf{Limited datasets.} While the datasets represent important clinical scenarios, they do not cover all EEG tasks or populations.
  \item \textbf{Preprocessing effects.} Normalization, filtering, and montage harmonization influence results and are held fixed in this study.
  \item \textbf{Hyperparameter sensitivity.} The choice of classifier head, regularization, and learning rate can affect the ranking of models.
\end{enumerate}

Despite these limitations, the benchmark provides actionable insights for practitioners and researchers developing EEG-based diagnostic systems.

% ==================================================================
\section{Conclusion}
% ==================================================================

This paper benchmarks state-of-the-art EEG foundation models on clinically challenging datasets, emphasizing robustness under class imbalance and distribution shift. Using controlled evaluation protocols and standardized preprocessing, we show that:

\begin{enumerate}
  \item Model performance varies substantially across tasks and is strongly shaped by prevalence, dataset structure, and task separability.
  \item EEGDM emerges as the most consistent performer, achieving strong results on TUH Seizure (0.874 AUROC), Siena (0.915 AUROC), and Bonn (0.994 AUROC).
  \item Rare-event detection (Siena 0.7\% prevalence) remains fundamentally challenging, with even the best model achieving minority-class F1 of only 0.11.
  \item Prevalence matching substantially affects calibration and ranking; AUROC is insufficient and precision-recall metrics must be prioritized.
  \item No model is universally dominant, and real clinical utility depends on careful threshold selection and realistic prevalence-aware evaluation.
\end{enumerate}

These findings emphasize that EEG benchmarking must move beyond isolated accuracy reports and report prevalence-aware metrics, calibration, and minority-class sensitivity. Real clinical utility depends on the ability to detect rare but consequential events with controlled false alarms, not simply on balanced accuracy in an artificially favorable benchmark.

This benchmark provides a rigorous and reproducible foundation for future evaluation of EEG foundation models. We recommend that future studies:
\begin{itemize}
  \item Report results on at least two datasets with different prevalence regimes.
  \item Include both precision-recall and ROC curves with confidence intervals.
  \item Explicitly discuss threshold selection and clinical operating points.
  \item Provide code and splits for reproducibility.
\end{itemize}

By establishing these standards, the EEG and broader biomedical signal processing community can move toward more clinically grounded model development and deployment.

% ==================================================================
\section*{Acknowledgments}
% ==================================================================
The authors thank the TUH, TUAB, Siena, and Bonn dataset providers for making these datasets available for research. This work was supported by \TODO{funding information}.

% ==================================================================
% References
% ==================================================================
\begin{thebibliography}{99}

\bibitem{obeid2016tuh}
I. Obeid and J. A. Picone, ``The temple university hospital EEG data corpus,'' \textit{Front. Neurosci.}, vol. 10, p. 196, 2016.

\bibitem{yang2024biot}
C. Yang, R. Westman, and A. Rousseau, ``BIOT: Biosignal transformer for cross-dataset EEG foundation models,'' \textit{arXiv}, 2024.

\bibitem{jiang2024labram}
W. Jiang, H. Feng, B. Ju, and T. Liu, ``LaBraM: Towards large brain models with vector-quantized EEG,'' \textit{arXiv}, 2024.

\bibitem{wang2024cbramod}
X. Wang, Z. Zhang, and Y. Wang, ``CBraMod: Criss-cross transformer for brain modelling in EEG,'' \textit{arXiv}, 2024.

\bibitem{helwan2025diffeeg}
A. Helwan et al., ``DiffEEG: A denoising-diffusion foundation model for transferable EEG representations,'' \textit{IEEE Trans. Biomed. Eng.}, 2025. [Online]. Available: \url{https://github.com/AbdulkaderHelwan/Diff-EEG-V2}

\bibitem{davis2006relationship}
J. Davis and M. Goadrich, ``The relationship between precision-recall and ROC curves,'' in \textit{Proc. 23rd Int. Conf. Mach. Learn.}, 2006, pp. 233--240.

\bibitem{saito2015precision}
T. Saito and M. Rehmsmeier, ``The precision-recall plot is more informative than the ROC plot when evaluating binary classifiers on imbalanced datasets,'' \textit{PLoS ONE}, vol. 10, no. 3, p. e0118432, 2015.

\bibitem{meier2017seizure}
R. Meier, H. Dittrich, A. Schulze-Bonhage, and A. Aertsen, ``Detecting epileptic seizures in long-term human EEG: A new approach to automatic online and real-time detection and classification of polymorphic seizure patterns,'' \textit{J. Clin. Neurophysiol.}, vol. 25, no. 2, pp. 119--131, 2017.

\bibitem{siena_dataset}
\TODO{Add formal Siena dataset citation},

\end{thebibliography}

\end{document}
